\section{Implementazione}
\label{chap:implementation}

L'applicazione è stata sviluppata adottando un'architettura client-server disaccoppiata (API RESTful e Single Page Application per il frontend) e utilizzando il linguaggio di programmazione
TypeScript sia per il frontend che per il backend. 
Per la parte di backend è stato sfruttato il
runtime BunJS con il framework Hono
per la creazione e la gestione del server e delle API
RESTful. Per la parte di frontend è stato utilizzato il
framework React (versione 19) con il bundler
Vite. Il database utilizzato è MongoDB,
gestito tramite l'ODM Mongoose.

L'acquisizione dei dati dai sensori IoT avviene tramite il
protocollo MQTT: il backend si sottoscrive al broker (topic \texttt{sensors/{sensorId}/readings}) e archivia le
rilevazioni ricevute in MongoDB, utilizzando una time-series
collection per lo storico delle misurazioni.

\subsection{Scelta del Tech Stack e Motivazioni Architetturali}
La scelta di usare TypeScript è scaturita dalla necessità di avere
un controllo maggiore sul type-checking e una maggiore manutenibilità
del codice.
Questo ci permette di catturare una serie di bug a compile-time, evitando errori di tipo durante l'esecuzione.\\

La scelta del framework Hono presenta alcune principali motivazioni: performance, portabilità e Developer Experience (DX).
Hono garantisce performance migliori rispetto a Express.js e risulta conforme agli standard Web API, facendo si che si possa fare l'hosting anche su piattaforme edge (come Clodflare Workers o Vercel functions).
Inoltre Hono presenta un client per il frontend tipizzato, cosi che le risposte delle API siano tipizzate end-to-end. Cosi facendo si ottene l'auto-completamento e convalida dei dei parametri e dei payload restituiti a compile-time, eliminando il rischio di disallineamenti tra i due layer.\\

Il runtime BunJS è stato scelto principalmente per l'ottima DX e per le sue perfomance. Bun esegue nativamente TypeScript senza aver bisogno di strumenti come ts-node o nodemon. Inoltre esso integra un package manager molto più prestante di npm.
Infine, l'inclusione di alcune API (come cron e hashing per le password) ha  ridotto la necessità di scaricare altre dipendenze esterne.\\

Per il frontend, è stato adottato React 19 data l'esigenza di realizzare una Single Page Application. Rispetto ad alternative come Vue.js, la preferenza per React è stata guidata dalla maggiore capillarità dell'ecosistema e dall'integrazione con librerie come TanStack Query, React Router e Shadcn UI (quest'ultima basata sulle primitive React). Questo ha reso la gestione di caching, UI e routing molto più semplice. Inoltre l'affiancamento a Vite come build tool consente di ottenere un HMR istantaneo, migliorando ulteriormente la DX.\\

La persistenza dei dati si basa su MongoDB e sull'ODM Mongoose. 
L'uso di MongoDB è stato è giustificato in particolare dalla presenza nativa delle Time-Series Collections, permettendo di gestire grandi volumi di dati di time-series, come le letture dei sensori. Ciò è risultato essenziale per il calcolo dell'efficienza di ogni edificio in maniera efficiente.
\\
Infine, l'acquisizione dei dati mediante protocollo MQTT risponde ai vincoli operativi delle architetture IoT. A differenza di HTTP, MQTT opera tramite un modello Publish/Subscribe su TCP rendendolo essenziale per la gestione dei dati provenienti dai sensori.

\subsection{\textit{Repository Organization}}
Il codice del progetto (\href{https://github.com/alterax05/wattguard-ing-sw}{github.com/alterax05/wattguard-ing-sw}) è organizzato secondo la seguente struttura:

\dirtree{%
  .1 /.
  .2 /apps.
  .3 /api\DTcomment{Backend Bun + Hono}.
  .4 build.ts\DTcomment{Script di build di produzione}.
  .4 /src.
  .5 index.ts\DTcomment{Entry point del server backend}.
  .5 /routes\DTcomment{Endpoint RESTful raggruppati per risorsa (\texttt{/api/v1})}.
  .5 /models\DTcomment{Modelli Mongoose}.
  .5 /services\DTcomment{Servizi di background (ingestione, simulatore, cron allarmi)}.
  .5 /middleware\DTcomment{Middleware (autenticazione JWT, ruoli, rate limit, i18n)}.
  .5 /lib\DTcomment{Logica di dominio (efficienza termica, report PDF/XLSX, CSV)}.
  .5 /email\DTcomment{Email via API HTTP (\texttt{Resend})}.
  .5 /auth\DTcomment{Utility autenticazione}.
  .5 /config\DTcomment{Configurazione OpenAPI 3.1 e variabili d'ambiente}.
  .5 /utils\DTcomment{Utility}.
  .5 /\_\_tests\_\_\DTcomment{Suite di test backend}.
  .3 /web\DTcomment{Frontend React + Vite}.
  .4 /src.
  .5 main.tsx\DTcomment{Entry point}.
  .5 App.tsx.
  .5 /pages\DTcomment{Pagine dell'interfaccia (Dashboard, Edifici, Sensori, \dots)}.
  .5 /components\DTcomment{Componenti UI di (\texttt{Shadcn UI})}.
  .5 /hooks\DTcomment{Hook custom e data fetching con (\texttt{TanStack Query})}.
  .5 /lib\DTcomment{Client Hono tipizzato e egestione errori}.
  .5 /styles\DTcomment{Fogli di stile globali (\texttt{Tailwind CSS})}.
  .4 vite.config.ts.
  .2 /shared\DTcomment{Package condiviso}.
  .3 /assets\DTcomment{Asset grafici (loghi)}.
  .3 /locales\DTcomment{Cataloghi di traduzione (\texttt{en}, \texttt{it}, \texttt{de})}.
  .3 /src.
  .4 index.ts.
  .4 error-codes.ts.
  .4 i18n.ts\DTcomment{Gestione della lingua e localizzazione}.
  .4 /schemas\DTcomment{Schemi di validazione \texttt{Zod} condivisi tra backend e frontend}.
  .2 /scripts\DTcomment{Script ausiliari}.
  .2 /docs\DTcomment{Deliverables e OpenAPI doc esportato}.
  .2 docker-compose.yml\DTcomment{Broker \texttt{MQTT} Mosquitto per lo sviluppo locale}.
  .2 .oxlintrc.json.
  .2 package.json.
  .2 tsconfig.base.json.
}

\subsection{Branching Strategy e organizzazione del lavoro}
    Per la gestione e condivisione del codice si è scelto di utilizzare
    \texttt{GitHub} e di adottare una strategia di branching basata
    su \textit{GitFlow}. In particolare si è deciso di utilizzare i seguenti
    \textit{branch}:
    \begin{description}
        \item[\texttt{main}] \textit{Branch} principale, contiene il codice
        stabile;
        \item[\texttt{develop}] \textit{Branch} di integrazione, sul quale
        confluiscono tutte le funzionalità sviluppate;
        \item[\texttt{feature/<nome>}] \textit{Branch} per lo sviluppo di
        nuove funzionalità, creati da \texttt{develop}
        \item[\texttt{fix/<nome>}] \textit{Branch} per le correzioni
        di errori
        \item[\texttt{docs/<nome>}] \textit{Branch} per la documentazione.
    \end{description}
    I commit seguono la specifica Conventional Commits
    (\texttt{feat:}, \texttt{fix:}, \texttt{chore:}, \dots).
    Ad ogni push sul ramo \texttt{develop} e sul ramo \texttt{main} vengono eseguite verifiche sulla qualità del codice tra cui controlli statici e test.

    Per la divisione delle aree di lavoro abbiamo individuato un responsabile principale mentre per le slide del pitch ogniuno ha contribuito in modo equo.
    \begin{itemize}
        \item \textbf{Michele Porcellato}: responsabile parte pratica del backend. Responsabile principale della scrittura del documento D2 e addetto alla prima scrittura dei requisiti funzionali e non funzionali per il deliverable D1.
        \item \textbf{Marco Piovesan}: Responsabile principale per il documento D1 e frontend. Inoltre è corresponsabile per documento D3. 
        \item \textbf{Alessio Pesaresi}: Responsabile principale per il documento D3, addetto alla stesura del business plan e alla realizzazione del video. Inoltre è corresponsabile per lo sviluppo del frontend.
    \end{itemize}
    
\subsection{\textit{Dependency}}
    \paragraph{Dipendenze backend e condivise} Per il \textit{backend} e il pacchetto condiviso sono state impiegati i seguenti moduli esterni:
    \begin{description}
        \item[\texttt{hono}] \textit{framework} web per la gestione delle API;
        \item[\texttt{hono-openapi}] per la generazione automatica della specifica \texttt{OpenAPI 3.1} partendo dalle definizioni delle rotte;
        \item[\texttt{@scalar/hono-api-reference}] per la documentazione interattiva (Scalar UI) esposta all'indirizzo \texttt{/api/v1/docs};
        \item[\texttt{hono-rate-limiter}] per la protezione da attacchi di brute-force (autenticazione, recupero password, inviti);
        \item[\texttt{mongoose}] libreria ODM (Object Data Modeling) per l'interazione con il database \texttt{MongoDB};
        \item[\texttt{mqtt}] client MQTT per la connessione al broker e l'acquisizione delle letture dei sensori \texttt{IoT};
        \item[\texttt{zod}] libreria per la validazione dei dati in input;
        \item[\texttt{resend}] per l'invio di email tramite API HTTP;
        \item[\texttt{google-auth-library}] per supporto del login con Google;
        \item[\texttt{exceljs} e \texttt{pdfkit}] per la generazione dei report di efficienza energetica in formato .xlsx e .pdf;
        \item[\texttt{i18next}] per il supporto multilingua sia sul frontend sia sul backend;
    \end{description}
    \paragraph{Dipendenze frontend} Per il \textit{frontend} le librerie principali includono:
    \begin{description}
        \item[\texttt{react}, \texttt{react-dom}] libreria frontend React;
        \item[\texttt{vite}] bundler e dev server;
        \item[\texttt{react-router-dom}] per la gestione del routing client-side;
        \item[\texttt{@tanstack/react-query}] per la gestione di fetching, caching, e sicronizzazione di dati con il backend;
        \item[\texttt{tailwindcss} (v4) e \texttt{radix-ui}] per la gestione dello stile dell'applicazione;
        \item[\texttt{recharts}] per la visualizzazione di grafici;
        \item[\texttt{leaflet} e \texttt{react-leaflet}] per la mappa degli edifici monitorati;
        \item[\texttt{sonner} e \texttt{next-themes}] per il sistema di notifiche \textit{toast} e la gestione del tema chiaro e scuro;
        \item[\texttt{react-i18next} e \texttt{i18next-browser-languagedetector}] per l'internazionalizzazione dell'interfaccia;
        \item[\texttt{react-hook-form} e \texttt{@hookform/resolvers}] per la gestione dei moduli presenti nel frontend e la validazione integrata tramite Zod;
        \item[\texttt{date-fns}] per la gestione di date e intervalli temporali;
        \item[\texttt{lucide-react}] catalogo completo di icone vettoriali per l'interfaccia utente;
    \end{description}
    \paragraph{Dipendenze di sviluppo}
    \begin{description}
        \item[\texttt{typescript}] per il controllo statico dei tipi;
        \item[\texttt{oxlint}] linter per l'analisi statica del codice;
        \item[\texttt{mongodb-memory-server}] per l'avvio di un'istanza MongoDB effimera durante l'esecuzione dei test di integrazione;
    \end{description}

\subsection{\textit{Database}}
    Il database è stato progettato per contenere le seguenti
    collezioni su MongoDB, gestite tramite Mongoose:
    \subsubsection{Utente}
        La collezione \texttt{users} viene usata per memorizzare i dati
        relativi agli utenti del sistema. Ogni documento della collezione
        contiene i seguenti campi:
        \begin{description}
            \item[\texttt{\_id}] (\texttt{ObjectId}) identificativo univoco dell'utente;
            \item[\texttt{email}] (\texttt{String}) email dell'utente, univoca;
            \item[\texttt{name}] (\texttt{String}) nome visualizzato dell'utente (facoltativo);
            \item[\texttt{role}] (\texttt{String, enum \texttt{admin}$|$\texttt{operator}}) ruolo dell'utente nel sistema;
            \item[\texttt{isDisabled}] (\texttt{Boolean}) \textit{flag} che indica se l'account è disabilitato;
            \item[\texttt{language}] (\texttt{String, enum \texttt{en}$|$\texttt{it}$|$\texttt{de}}) preferenza linguistica dell'utente (facoltativa);
            \item[\texttt{passwordHash}] (\texttt{String}) hash \texttt{bcrypt} della password;
            \item[\texttt{googleSub}] (\texttt{String}) identificativo univoco Google
            \item[\texttt{lastLoginAt}] (\texttt{Date}) timestamp dell'ultimo accesso effettuato;
            \item[\texttt{passwordUpdatedAt}] (\texttt{Date}) data dell'ultimo cambio password;
            \item[\texttt{createdAt} / \texttt{updatedAt}] (\texttt{Date}) timestamp di creazione e modifica;
        \end{description}
    \subsubsection{Invito}
        La collezione \texttt{invites} viene usata per memorizzare gli inviti
        inviati via email per la registrazione dei nuovi utenti. Ogni documento
        contiene i seguenti campi:
        \begin{description}
            \item[\texttt{\_id}] (\texttt{ObjectId}) identificativo univoco dell'invito;
            \item[\texttt{email}] (\texttt{String}) email dell'invitato
            \item[\texttt{role}] (\texttt{String, enum \texttt{admin}$|$\texttt{operator}}) ruolo assegnato all'utente al momento dell'accettazione;
            \item[\texttt{tokenHash}] (\texttt{String}) hash \texttt{SHA-256} del token segreto generato casualmente;
            \item[\texttt{status}] (\texttt{String, enum \texttt{pending}$|$\texttt{accepted}$|$\texttt{revoked}$|$\texttt{expired}}) stato corrente dell'invito;
            \item[\texttt{expiresAt}] (\texttt{Date}) scadenza dell'invito (7 giorni dalla creazione), associata a un indice \textbf{TTL} per la cancellazione automatica;
            \item[\texttt{createdBy}] (\texttt{ObjectId}) amministratore che ha generato l'invito (riferimento a \texttt{User});
            \item[\texttt{acceptedAt}] (\texttt{Date}) data e ora di avvenuta accettazione dell'invito;
            \item[\texttt{createdAt} / \texttt{updatedAt}] (\texttt{Date}) timestamp di creazione e modifica;
        \end{description}
    \subsubsection{TokenResetPassword}
        La collezione \texttt{passwordresettokens} viene usata per memorizzare
        i token monouso utilizzati durante il reset della password. Ogni documento
        contiene i seguenti campi:
        \begin{description}
            \item[\texttt{\_id}] (\texttt{ObjectId}) identificativo univoco del token;
            \item[\texttt{userId}] (\texttt{ObjectId}) utente a cui il token è associato (riferimento a \texttt{User});
            \item[\texttt{tokenHash}] (\texttt{String}) hash \texttt{SHA-256} del token monouso, univoco;
            \item[\texttt{expiresAt}] (\texttt{Date}) scadenza del token (1 ora dalla richiesta), associata a un indice \textbf{TTL} per la rimozione automatica;
            \item[\texttt{usedAt}] (\texttt{Date}) data e ora di effettivo utilizzo, al fine di prevenire attacchi di riuso del token;
            \item[\texttt{createdAt} / \texttt{updatedAt}] (\texttt{Date}) timestamp di creazione e modifica;
        \end{description}
    \subsubsection{Edificio}
        La collezione \texttt{buildings} viene usata per memorizzare
        l'anagrafica, i parametri termici e la localizzazione degli edifici monitorati. Ogni documento contiene i seguenti campi:
        \begin{description}
            \item[\texttt{\_id}] (\texttt{ObjectId}) identificativo univoco dell'edificio;
            \item[\texttt{name}] (\texttt{String}) nome identificativo dell'edificio;
            \item[\texttt{address}] (\texttt{String}) indirizzo stradale completo;
            \item[\texttt{surface}] (\texttt{Number}) superficie utile riscaldata in \(m^2\);
            \item[\texttt{ceilingHeight}] (\texttt{Number}) altezza media dei soffitti in metri (minimo 0.5, massimo 20, default 3.0);
            \item[\texttt{location}] (\texttt{Object}) posizione geografica \texttt{GeoJSON} \texttt{\{type: 'Point', coordinates: [lng, lat]\}}, indicizzata con indice geospaziale sferico \texttt{2dsphere};
            \item[\texttt{buildingType}] (\texttt{ObjectId}) riferimento alla tipologia di edificio (riferimento a \texttt{BuildingType});
            \item[\texttt{heatingSystemType}] (\texttt{String}) tipo di impianto termico installato (es. pompa di calore, teleriscaldamento, caldaia a gas);
            \item[\texttt{constructionYear}] (\texttt{Number}) anno di costruzione dell'immobile;
            \item[\texttt{geographicZone}] (\texttt{String}) zona geografica o climatica dell'edificio;
            \item[\texttt{status}] (\texttt{String, enum \texttt{active}$|$\texttt{inactive}$|$\texttt{decommissioned}}) stato operativo dell'edificio;
            \item[\texttt{efficiencyThresholds}] (\texttt{Object}) configurazione dei limiti di efficienza energetica \texttt{\{enabled: Boolean, minCop: Number\}};
            \item[\texttt{createdBy} / \texttt{updatedBy}] (\texttt{ObjectId}) riferimenti agli utenti che hanno creato e aggiornato l'edificio;
            \item[\texttt{createdAt} / \texttt{updatedAt}] (\texttt{Date}) timestamp di creazione e modifica;
        \end{description}.
    \subsubsection{TipologiaEdificio}
        La collezione \texttt{buildingtypes} viene usata per memorizzare le
        categorie di edifici per la ricerca. Ogni documento contiene i seguenti campi:
        \begin{description}
            \item[\texttt{\_id}] (\texttt{ObjectId}) identificativo univoco della tipologia;
            \item[\texttt{name}] (\texttt{String}) nome della tipologia, univoco nel sistema;
            \item[\texttt{description}] (\texttt{String}) descrizione delle caratteristiche della tipologia (facoltativa);
            \item[\texttt{createdAt} / \texttt{updatedAt}] (\texttt{Date}) timestamp di creazione e modifica;
        \end{description}
    \subsubsection{Sensore}
        La collezione \texttt{sensors} viene usata per memorizzare i sensori
        installati negli edifici, le eventuali soglie di allerta e lo stato di operatività. Ogni documento contiene i seguenti campi:
        \begin{description}
            \item[\texttt{\_id}] (\texttt{ObjectId}) identificativo univoco del sensore;
            \item[\texttt{building}] (\texttt{ObjectId}) riferimento all'edificio in cui il sensore è installato (indicizzato, ref \texttt{Building});
            \item[\texttt{sensorType}] (\texttt{String, enum \texttt{internal\_temp}$|$\texttt{external\_temp}$|$\texttt{energy\_meter}$|$\texttt{gas\_meter}}) tipo di sensore supportato dall'applicazione;
            \item[\texttt{location}] (\texttt{String}) collocazione fisica del sensore nell'edificio;
            \item[\texttt{serialNumber}] (\texttt{String}) numero di serie hardware del sensore;
            \item[\texttt{installationDate}] (\texttt{Date}) data di installazione del dispositivo (default \texttt{Date.now});
            \item[\texttt{status}] (\texttt{String, enum \texttt{active}$|$\texttt{inactive}$|$\texttt{maintenance}$|$\texttt{error}}) stato operativo del sensore;
            \item[\texttt{lastReading}] (\texttt{Object}) ultima rilevazione ricevuta \texttt{\{value, timestamp, unit\}}, denormalizzata per la lettura istantanea senza interrogare la collezione \texttt{sensorreadings};
            \item[\texttt{transmissionInterval}] (\texttt{Number}) intervallo nominale di trasmissione in secondi;
            \item[\texttt{minThreshold} / \texttt{maxThreshold}] (\texttt{Number}) soglie minima e massima per il rilevamento di anomalie e la generazione di allarmi;
            \item[\texttt{createdBy} / \texttt{updatedBy}] (\texttt{ObjectId}) utenti responsabili della registrazione e modifica del dispositivo;
            \item[\texttt{createdAt} / \texttt{updatedAt}] (\texttt{Date}) timestamp di creazione e aggiornamento;
        \end{description}
    \subsubsection{Letture sensori}
        La collezione \texttt{sensorreadings} è una \emph{time-series
        collection} nativa di MongoDB usata per memorizzare lo storico delle misurazioni dei
        sensori. Il contenuto di questa collezione ha scadenza automatica di 365 giorni (modificabile dalle impostazioni dell'applicazione). Ogni documento contiene:
        \begin{description}
            \item[\texttt{\_id}] (\texttt{ObjectId}) identificativo univoco della rilevazione;
            \item[\texttt{timestamp}] (\texttt{Date}) istante temporale esatto della misurazione in formato UTC (\texttt{timeField});
            \item[\texttt{value}] (\texttt{Number}) valore numerico rilevato;
            \item[\texttt{unit}] (\texttt{String}) unità di misura fisica (ad esempio \texttt{\textdegree C}, \texttt{kW}, \texttt{m$^3$});
            \item[\texttt{metadata}] (\texttt{Object}) metadati di contesto per il partizionamento efficiente \texttt{\{sensor (ObjectID), building (ObjectID), sensorType (String)\}} (\texttt{metaField});
        \end{description}
    \subsubsection{Allerte}
        La collezione \texttt{alerts} viene usata per memorizzare gli allarmi
        generati dal sistema a seguito di superamento soglie preimpostate o all'abbassamento dell'efficienza energetica. Ongi documento contiene:
        \begin{description}
            \item[\texttt{\_id}] (\texttt{ObjectId}) identificativo univoco dell'allarme;
            \item[\texttt{building}] (\texttt{ObjectId}) edificio a cui l'allarme si riferisce (riferimento a \texttt{Building});
            \item[\texttt{sensor}] (\texttt{ObjectId}) sensore che ha causato l'anomalia (riferimento a Sensor);
            \item[\texttt{type}] (\texttt{String}) tipologia dell'anomalia;
            \item[\texttt{thresholdType}] (\texttt{String, enum \texttt{min}$|$\texttt{max}}) soglia superata (per gli allarmi di soglia);
            \item[\texttt{severity}] (\texttt{String, enum \texttt{low}$|$\texttt{medium}$|$\texttt{high}$|$\texttt{critical}}) livello di gravità calcolato dinamicamente;
            \item[\texttt{sensorType}] (\texttt{String}) tipologia del sensore coinvolto (facoltativa);
            \item[\texttt{location}] (\texttt{String}) collocazione fisica del sensore al momento dell'evento (facoltativa);
            \item[\texttt{value}] (\texttt{Number}) valore riscontrato che ha innescato l'allarme;
            \item[\texttt{unit}] (\texttt{String}) unità di misura del valore anomalo;
            \item[\texttt{limit}] (\texttt{Number}) soglia di riferimento superata;
            \item[\texttt{status}] (\texttt{String, enum \texttt{active}$|$\texttt{acknowledged}$|$\texttt{resolved}}) stato di gestione dell'allarme;
            \item[\texttt{acknowledgedBy} / \texttt{acknowledgedAt}] operatore e data di presa in carico;
            \item[\texttt{resolvedBy} / \texttt{resolvedAt}] operatore e data di risoluzione;
            \item[\texttt{createdAt} / \texttt{updatedAt}] (\texttt{Date}) timestamp di creazione e aggiornamento;
        \end{description}
    \subsubsection{ConfigurazioneSistema}
        La collezione \texttt{systemconfigs} è un documento \emph{singleton}
        gestito tramite il metodo statico \texttt{SystemConfig.getOrCreate()}, che memorizza la configurazione globale del sistema. Il documento contiene:
        \begin{description}
            \item[\texttt{polling.intervalSeconds}] (\texttt{Number}) intervallo di aggiornamento automatico per il \textit{frontend} (default 90, range 10--3600);
            \item[\texttt{polling.autoPollingEnabled}] (\texttt{Boolean}) abilitazione del \textit{polling} automatico nella dashboard;
            \item[\texttt{notifications.emailEnabled}] (\texttt{Boolean}) abilitazione globale delle notifiche di allarme via email;
            \item[\texttt{database.dataRetentionDays}] (\texttt{Number}) periodo di conservazione dei delle misurazioni dei sensori (default 365 giorni);
            \item[\texttt{createdAt} / \texttt{updatedAt}] (\texttt{Date}) timestamp di audit;
        \end{description}

%\subsection{\textit{Acquisizione dei dati dai sensori}}
%    \hl{Utile?}
%    L'acquisizione della telemetria IoT si articola attraverso due canali complementari che convergono su una pipeline di elaborazione comune:
%    \begin{itemize}
%        \item \textbf{Broker MQTT}: in ambiente di produzione e di sviluppo (tramite container \texttt{eclipse-mosquitto:2.0} configurato in \texttt{docker-compose.yml}), il backend si connette al broker e si sottoscrive con wildcard al topic \texttt{sensors/+/readings}. La logica di ascolto e decodifica è gestita dal modulo \texttt{apps/api/src/lib/mqtt.ts};
%        \item \textbf{Simulatore In-Process}: per agevolare lo sviluppo locale e i test di carico senza richiedere un broker MQTT esterno, il backend integra un simulatore interno (\texttt{apps/api/src/services/simulator.ts}) attivabile tramite la variabile \texttt{SIMULATOR\_ENABLED=true}. Tale componente genera a intervalli realistici misurazioni fisiche coerenti per tutti gli edifici (temperature interne/esterne e consumi) ed effettua l'auto-seed dei dati di prova in ambiente di sviluppo.
%    \end{itemize}
%    Tutte le rilevazioni vengono convogliate nel servizio centralizzato di ingestione (\texttt{apps/api/src/services/reading-service.ts}). Per ogni lettura il sistema esegue:
%    \begin{enumerate}
%        \item La validazione e il recupero del sensore dal database;
%        \item La verifica del valore rispetto alle soglie operative \texttt{minThreshold} e \texttt{maxThreshold};
%        \item Qualora il valore risulti fuori soglia e non esista già un allarme attivo per la medesima anomalia (grazie all'indice composto di deduplicazione), il calcolo dinamico della severità tramite il modulo \texttt{alert-severity.ts} e la generazione di un nuovo \texttt{Alert};
%        \item Se le notifiche via email sono abilitate nella configurazione globale, l'invio immediato di un'email di allarme localizzata tramite il servizio \texttt{notification-service.ts} che si interfaccia con le API di \texttt{Resend};
%        \item L'archiviazione della misurazione nella time-series collection \texttt{sensorreadings};
%        \item L'aggiornamento denormalizzato di \texttt{lastReading} sul documento del sensore, ripristinandone automaticamente lo stato ad \texttt{active} qualora fosse stato precedentemente marcato come non in linea.
%    \end{enumerate}
%    In aggiunta al controllo reattivo sulle singole letture, il sistema impiega un task in background periodico (\texttt{apps/api/src/services/efficiency-alert-service.ts}), schedulato nativamente con \texttt{Bun.cron}, che valuta periodicamente il Coefficiente di Prestazione (COP) stimato per ciascun edificio rispetto alla soglia minima configurata in \texttt{efficiencyThresholds.minCop}, generando allarmi proattivi di tipo \texttt{efficiency\_cop\_low} in caso di calo del rendimento termico.

\subsection{Testing}
Le API e le logiche di dominio implementate presentano una test-suite che consente di verificarne il corretto funzionamento. I test sono utilizzati all'interno della configurazione di CI/CD

L'implementazione dei test è organizzata in file \texttt{.test.ts} all'interno della directory \newline
\texttt{apps/api/src/\_\_tests\_\_/} (strutturata specularmente al codice sorgente). 
Durante la fase di testing, tutti i servizi esterni, come login con Google, API Resend e broker MQTT, sono simulati (uso di mock) eccetto per il database. Per simulare il database si è ricorso all'uso della libreria \texttt{mongodb-memory-server}, avviando un processo MongoDB effimero per la durata dei test.

La suite conta complessivamente 149 casi di test distribuiti su 18 file, ottenendo oltre il 70\% di code coverage.

\begingroup
\scriptsize
\setlength{\tabcolsep}{2.5pt}
\renewcommand{\arraystretch}{1.2}

% Definizione pesi colonne: totale pesi X = 5.0 (esattamente pari a 5 colonne X)
\begin{xltabular}{\textwidth}{
  |p{0.6cm}
  |>{\hsize=1.2\hsize}X
  |>{\hsize=0.8\hsize}X
  |>{\hsize=1.1\hsize}X
  |>{\centering\arraybackslash}p{0.75cm}
  |>{\hsize=0.95\hsize}X
  |>{\hsize=0.95\hsize}X
  |p{1.35cm}|
}

    \hline
    \textbf{N.} & \textbf{Descrizione} & \textbf{\textit{Test Data}} & \textbf{Precondizioni} & \textbf{Dip.} & \textbf{Risultato Atteso} & \textbf{Risultato Riscontrato} & \textbf{Note} \\
    \hline
    \endfirsthead

    \multicolumn{8}{l}{\textit{Tabella continuata dalla pagina precedente}} \\
    \hline
    \textbf{N.} & \textbf{Descrizione} & \textbf{\textit{Test Data}} & \textbf{Precondizioni} & \textbf{Dip.} & \textbf{Risultato Atteso} & \textbf{Risultato Riscontrato} & \textbf{Note} \\
    \hline
    \endhead

    \hline
    \multicolumn{8}{|r|}{\textit{Tabella continuata nella pagina successiva}}\\
    \hline
    \endfoot

    \hline
    \endlastfoot

    1.1 & \texttt{POST /api/v1/auth/session} login con credenziali valide & \textit{email}, \textit{password} & Utente registrato con password corretta & \texttt{DB} & La richiesta ritorna status \texttt{200} mentre nel body è presente token JWT e dati del profilo& La richiesta ritorna status \texttt{200} mentre nel body è presente token JWT e dati del profilo& RF1 \\
    \hline
    1.2 & \texttt{POST /api/v1/auth/session} login con password errata & \textit{email}, \textit{password} & Utente esistente, password non corrispondente & \texttt{DB} & La richiesta ritorna status \texttt{401}  e come codice d'errore "invalid\_credentials" & La richiesta ritorna status \texttt{401}  e come codice d'errore "invalid\_credentials" & RF1 \\
    \hline
    1.3 & \texttt{POST /api/v1/auth/session} login con Google ID token & \texttt{idToken} Google & Token valido emesso da Google (GSI) e utente esistente & \texttt{DB}, \texttt{Google} & La richesta ritorna status \texttt{200}, utente autenticato e token JWT nel body & La richesta ritorna status \texttt{200}, utente autenticato e token JWT nel body & RF1 \\
    \hline
    1.4 & \texttt{POST /api/v1/auth/session} login Google utente inesistente & \textit{idToken} Google & Token valido ma nessuna email registrata & \texttt{DB}, \texttt{Google} & La richiesta ritorna status \texttt{404} e codice d'errore \texttt{oauth\_no\_account} & La richiesta ritorna status \texttt{404} e codice d'errore \texttt{oauth\_no\_account} & RF1 \\
    \hline
    1.5 & \texttt{POST /api/v1/auth/session} login Google account disabilitato & \textit{idToken} Google & Token valido, utente con \texttt{isDisabled=true} & \texttt{DB}, \texttt{Google} & La richiesta ritorna status \texttt{403} con errore \texttt{oauth\_account\_disabled} &La richiesta ritorna status \texttt{403} con errore \texttt{oauth\_account\_disabled} & RF11 \\
    \hline
    1.6 & \texttt{GET /api/v1/auth/session} profilo con token valido & \textit{JWT token} & Token valido, utente attivo & \texttt{DB} & La richiesta ritorna status \texttt{200} e i dati del profilo & La richiesta ritorna status \texttt{200} e i dati del profilo & RF1 \\
    \hline
    1.7 & \texttt{GET /api/v1/auth/session} profilo senza token & - & Nessun token nell'header \texttt{Authorization} & - & La richiesta ritorna status \texttt{401} "Unauthorized" & La richiesta ritorna status \texttt{401} "Unauthorized" & RF1 \\
    \hline
    1.8 & \texttt{PATCH /api/v1/auth/session} modifica profilo & \textit{name}, \textit{language} & Utente autenticato & \texttt{DB} & La richiesta ritorna status \texttt{200} e il profilo aggiornato mentre status \texttt{400} con dati non validi & La richiesta ritorna status \texttt{200} e il profilo aggiornato mentre status \texttt{400} con dati non validi & RF0 \\
    \hline
    1.9 & \texttt{GET \nolinkurl{/api/v1/auth/google/config}} configurazione federata & - & - & - & La risposta ritorna status \texttt{200} e il client ID Google configurato & La risposta ritorna status \texttt{200} e il client ID Google configurato  & RF1 \\
    \hline
    1.10 & \texttt{GET /api/v1/invites?token=} validazione invito & \textit{token} & Invito \texttt{pending} non scaduto & \texttt{DB} & La richiesta ritorna status \texttt{200} e dati invito con email & La richiesta ritorna status \texttt{200} e dati invito con email & RF11 \\
    \hline
    1.11 & \texttt{PATCH /api/v1/invites/:id} setup credenziali locali & \textit{token}, \textit{password}, \textit{name} & Invito valido \texttt{pending} & \texttt{DB} & La richiesta ritorna status \texttt{200}, l'account attivato e token JWT nel body & La richiesta ritorna status \texttt{200}, l'account attivato e token JWT nel body & RF11 \\
    \hline
    1.12 & \texttt{PATCH /api/v1/invites/:id} accettazione invito tramite Google & \textit{idToken} Google & Invito valido in attesa di registrazione e nessun account preesistente & \texttt{DB}, \texttt{Google} & La richiesta ritorna status \texttt{200}, utente collegato e token JWT nel body & La richiesta ritorna status \texttt{200}, utente collegato e token JWT nel body & RF11 \\
    \hline
    2.1 & \texttt{GET /api/v1/users} elenco gli utenti (admin) & \textit{JWT token} & Utente \texttt{admin} autenticato & \texttt{DB} & La risposta ritorna status \texttt{200} e lista degli utenti priva degli hash della password & La risposta ritorna status \texttt{200} e lista degli utenti priva degli hash della password & RF11 \\
    \hline
    2.2 & \texttt{GET /api/v1/users/:id} dettaglio utente & \texttt{id} & Utente esistente, ruolo \texttt{admin} & \texttt{DB} & La richiesta ritorna status \texttt{200} assieme all'utente con tutti i suoi dettagli & La richiesta ritorna status \texttt{200} assieme all'utente con tutti i suoi dettagli & RF11 \\
    \hline
    2.3 & \texttt{GET /api/v1/users/:id} utente inesistente & \texttt{id} & ID non presente nel DB & \texttt{DB} & La richesta ritorna status \texttt{404} e il codice d'errore \texttt{user\_not\_found} & La richesta ritorna status \texttt{404} e il codice d'errore \texttt{user\_not\_found} & RF11 \\
    \hline
    2.4 & \texttt{PATCH /api/v1/users/:id} modifica ruolo & \texttt{id}, \textit{role} & Utente esistente diverso dall'admin corrente & \texttt{DB} & La richiesta ritorna status \texttt{200} e il ruolo aggiornato & La richiesta ritorna status \texttt{200} e il ruolo aggiornato & RF11 \\
    \hline
    2.5 & \texttt{PATCH /api/v1/users/:id} disabilitazione utente & \texttt{id}, \textit{isDisabled} & Utente esistente diverso dall'admin corrente & \texttt{DB} & La risposta ritorna status \texttt{200} e l'utente disabilitato & La risposta ritorna status \texttt{200} e l'utente disabilitato & RF11 \\
    \hline
    2.6 & \texttt{PATCH /api/v1/users/:id} modifica del proprio account & \texttt{id}, \textit{role} & L'\texttt{id} coincide con l'admin richiedente & \texttt{DB} & La richiesta ritorna status \texttt{422} con errore \texttt{cannot\_update\_own\_account} & La richiesta ritorna status \texttt{422} con errore \texttt{cannot\_update\_own\_account} & RF11 \\
    \hline
    2.7 & \texttt{PATCH /api/v1/users/:id} utente inesistente & \texttt{id} & ID non presente nel DB & \texttt{DB} & La richiesta ritorna status \texttt{404} e errore \texttt{user\_not\_found} & La richiesta ritorna status \texttt{404} e errore \texttt{user\_not\_found} & RF11 \\
    \hline
    2.8 & \texttt{PATCH /api/v1/users/:id} ruolo non valido & \textit{role} & Ruolo fuori enum \texttt{admin}$|$\texttt{operator} & \texttt{DB} & La risposta ritorna status \texttt{400} & La risposta ritorna status \texttt{400} & RF11 \\
    \hline
    2.9 & \texttt{POST /api/v1/auth/session} login utente disabilitato & \textit{email}, \textit{password} & Utente con \texttt{isDisabled=true} & \texttt{DB} & La richiesta ritorna status \texttt{403} & La richiesta ritorna status \texttt{403} & RF11 \\
    \hline
    2.10 & \texttt{DELETE /api/v1/users/:id} cancellazione utente & \texttt{id} & Utente esistente diverso dall'admin corrente & \texttt{DB} & La richiesta ritorna soltanto con status \texttt{204} & La richiesta ritorna soltanto con status \texttt{204} & RF11 \\
    \hline
    2.11 & \texttt{DELETE /api/v1/users/:id} cancellazione di s\'e stesso & \texttt{id} & L'\texttt{id} coincide con l'admin corrente & \texttt{DB} & La richiesta ritorna con status \texttt{422} con codice d'errore \texttt{cannot\_delete\_own\_account} & La richiesta ritorna con status \texttt{422} con codice d'errore \texttt{cannot\_delete\_own\_account} & RF11 \\
    \hline
    2.12 & \texttt{DELETE /api/v1/users/:id} utente inesistente & \texttt{id} & ID non presente nel DB & \texttt{DB} & La richiesta ritorna con status \texttt{404} e codice di errore \texttt{user\_not\_found} & La richiesta ritorna con status \texttt{404} e codice di errore \texttt{user\_not\_found} & RF11 \\
    \hline
    2.13 & \texttt{POST /api/v1/invites} creazione invito & \textit{email}, \textit{role} & Email non registrata, nessun invito pendente & \texttt{DB}, \texttt{Resend} & La richiesta ritorna con status \texttt{201}, creando con successo l'invito e inviando la mail & La richiesta ritorna con status \texttt{201}, creando con successo l'invito e inviando la mail & RF11 \\
    \hline
    2.14 & \texttt{POST /api/v1/invites} email gi\'a registrata & \textit{email}, \textit{role} & Email di utente esistente & \texttt{DB} & La richiesta ritorna status \texttt{409} e errore \texttt{user\_email\_exists} & La richiesta ritorna status \texttt{409} e errore \texttt{user\_email\_exists} & RF11 \\
    \hline
    2.15 & \texttt{POST /api/v1/invites} invito pendente esistente & \textit{email}, \textit{role} & Invito \texttt{pending} per la stessa email & \texttt{DB} & La richiesta ritorna status \texttt{409} con codice d'errore \texttt{invite\_pending\_exists} & La richiesta ritorna status \texttt{409} con codice d'errore \texttt{invite\_pending\_exists} & RF11 \\
    \hline
    2.16 & \texttt{POST /api/v1/invites} email non valida & \textit{email} & Formato email errato & \texttt{DB} & La richiesta ritorna con status \texttt{400} dovuto ad un errore di validazione & La richiesta ritorna con status \texttt{400} dovuto ad un errore di validazione & RF11, RNF5 \\
    \hline
    2.17 & \texttt{POST /api/v1/invites} ruolo non valido & \textit{role} & Ruolo fuori enum & \texttt{DB} & La richiesta ritorna con status \texttt{400} dovuto ad un errore di validazione & La richiesta ritorna con status \texttt{400} dovuto ad un errore di validazione & RF11, RNF5 \\
    \hline
    2.18 & \texttt{POST /api/v1/invites} email mancante & - & Campo \textit{email} assente & \texttt{DB} & La richiesta ritorna con status \texttt{400} dovuto ad un errore di validazione & La richiesta ritorna con status \texttt{400} dovuto ad un errore di validazione & RF11, RNF5 \\
    \hline
    2.19 & \texttt{POST /api/v1/invites} ruolo mancante & - & Campo \textit{role} assente & \texttt{DB} & La richiesta ritorna con status \texttt{400} dovuto ad un errore di validazione & La richiesta ritorna con status \texttt{400} dovuto ad un errore di validazione & RF11, RNF5 \\
    \hline
    2.20 & \texttt{GET /api/v1/invites} elenco inviti & - & Ruolo \texttt{admin} & \texttt{DB} & La richiesta ritorna status \texttt{200} e l'elenco completo degli inviti nel body  & La richiesta ritorna status \texttt{200} e l'elenco completo degli inviti nel body  & RF11 \\
    \hline
    2.21 & \texttt{GET /api/v1/invites/:id} dettaglio invito & \texttt{id} & Invito esistente, ruolo \texttt{admin} & \texttt{DB} & La richiesta ritorna status \texttt{200}, e nel body i dettagli dell'invito & La richiesta ritorna status \texttt{200}, e nel body i dettagli dell'invito & RF11 \\
    \hline
    2.22 & \texttt{GET /api/v1/invites/:id} invito inesistente & \texttt{id} & ID non presente nel DB & \texttt{DB} & La richiesta ritorna status \texttt{404} e codice d'errore \texttt{invite\_not\_found} & La richiesta ritorna status \texttt{404} e codice d'errore \texttt{invite\_not\_found} & RF11 \\
    \hline
    2.23 & \texttt{DELETE /api/v1/invites/:id} revoca invito gi\'a revocato & \texttt{id} & Invito in stato \texttt{revoked} & \texttt{DB} & Errore \texttt{409} transizione vietata & Errore \texttt{409} transizione vietata & RF11  \\
    \hline
    2.24 & \texttt{DELETE /api/v1/invites/:id} revoca invito pendente & \texttt{id} & Invito in stato \texttt{pending} & \texttt{DB} & La richiesta ritorna lo status \texttt{200}, stato \texttt{revoked} & La richiesta ritorna lo status \texttt{200}, stato \texttt{revoked} & RF11 \\
    \hline
    2.25 & \texttt{DELETE /api/v1/invites/:id} invito inesistente & \texttt{id} & ID non presente nel DB & \texttt{DB} & Errore \texttt{404} invito non trovato & Errore \texttt{404} invito non trovato & RF11 \\
    \hline
    2.26 & \texttt{GET /api/v1/invites?token=} lookup token valido & \textit{token} & Token grezzo corrispondente a invito \texttt{pending} & \texttt{DB} & La richiesta ritorna lo status \texttt{200}, dati invito & La richiesta ritorna lo status \texttt{200}, dati invito & RF11 \\
    \hline
    2.27 & \texttt{GET /api/v1/invites?token=} token sconosciuto & \textit{token} & Nessun invito corrispondente & \texttt{DB} & Errore \texttt{404} token non valido & Errore \texttt{404} token non valido & RF11 \\
    \hline
    2.28 & \texttt{PATCH /api/v1/invites/:id} token di altro invito & \textit{token} & Token legato a invito diverso & \texttt{DB} & Errore \texttt{404} token non valido & Errore \texttt{404} token non valido & RF11 \\
    \hline
    3.1 & \texttt{GET \nolinkurl{/api/v1/building-types}} lista tipologie & - & Ruolo \texttt{admin} & \texttt{DB} & La richiesta ritorna lo status \texttt{200}, array tipologie & La richiesta ritorna lo status \texttt{200}, array tipologie & RF4 \\
    \hline
    3.2 & \texttt{GET \nolinkurl{/api/v1/building-types/:id}} tipologia inesistente & \texttt{id} & ID non presente nel DB & \texttt{DB} & Errore \texttt{404} non trovata & Errore \texttt{404} non trovata & RF5 \\
    \hline
    3.3 & \texttt{GET \nolinkurl{/api/v1/building-types/:id}} ObjectId malformato & \texttt{id} & ID non valido & - & La richiesta ritorna con status \texttt{400} dovuto ad un errore di validazione & La richiesta ritorna con status \texttt{400} dovuto ad un errore di validazione & RF5, RNF5 \\
    \hline
    3.4 & \texttt{POST \nolinkurl{/api/v1/building-types}} creazione tipologia & \textit{name}, \textit{desc.} & Ruolo \texttt{admin}, nome univoco & \texttt{DB} & La richiesta ritorna lo status \texttt{201}, tipologia creata & La richiesta ritorna lo status \texttt{201}, tipologia creata & RF5 \\
    \hline
    3.5 & \texttt{POST \nolinkurl{/api/v1/building-types}} nome duplicato & \textit{name} & Nome gi\'a esistente & \texttt{DB} & Errore \texttt{409} violazione unicit\'a & Errore \texttt{409} violazione unicit\'a & RF5 \\
    \hline
    3.6 & \texttt{POST \nolinkurl{/api/v1/building-types}} dati non validi & \textit{name} & Nome vuoto o assente & \texttt{DB} & La richiesta ritorna con status \texttt{400} dovuto ad un errore di validazione & La richiesta ritorna con status \texttt{400} dovuto ad un errore di validazione & RF5, RNF5 \\
    \hline
    3.7 & \texttt{PATCH \nolinkurl{/api/v1/building-types/:id}} modifica tipologia & \texttt{id}, \textit{name} & Tipologia esistente, ruolo \texttt{admin} & \texttt{DB} & La richiesta ritorna lo status \texttt{200}, tipologia aggiornata & La richiesta ritorna lo status \texttt{200}, tipologia aggiornata & RF5 \\
    \hline
    3.8 & \texttt{PATCH \nolinkurl{/api/v1/building-types/:id}} nome duplicato & \textit{name} & Nome gi\'a in uso & \texttt{DB} & Errore \texttt{409} violazione unicit\'a & Errore \texttt{409} violazione unicit\'a & RF5 \\
    \hline
    3.9 & \texttt{PATCH \nolinkurl{/api/v1/building-types/:id}} tipologia inesistente & \texttt{id} & ID non presente nel DB & \texttt{DB} & Errore \texttt{404} non trovata & Errore \texttt{404} non trovata & RF5 \\
    \hline
    3.10 & \texttt{DELETE \nolinkurl{/api/v1/building-types/:id}} cancellazione tipologia & \texttt{id} & Tipologia esistente non referenziata & \texttt{DB} & La richiesta ritorna lo status \texttt{200}, tipologia rimossa & La richiesta ritorna lo status \texttt{200}, tipologia rimossa & RF5 \\
    \hline
    3.11 & \texttt{DELETE \nolinkurl{/api/v1/building-types/:id}} tipologia in uso & \texttt{id} & Referenziata da almeno un edificio & \texttt{DB} & Errore \texttt{400} vincolo referenziale & Errore \texttt{400} vincolo referenziale & RF5 \\
    \hline
    3.12 & \texttt{DELETE \nolinkurl{/api/v1/building-types/:id}} tipologia inesistente & \texttt{id} & ID non presente nel DB & \texttt{DB} & Errore \texttt{404} non trovata & Errore \texttt{404} non trovata & RF5 \\
    \hline
    4.1 & \texttt{POST /api/v1/buildings} creazione edificio & \textit{name}, \textit{addr}, \textit{surf}, \textit{loc}, \textit{type} & Dati validi, GeoJSON valido, tipologia esistente, ruolo \texttt{admin} & \texttt{DB} & La richiesta ritorna lo status \texttt{201}, edificio creato & La richiesta ritorna lo status \texttt{201}, edificio creato & RF5 \\
    \hline
    4.2 & \texttt{POST /api/v1/buildings} creazione con ruolo operator & \textit{name}, \textit{addr}, \textit{surf}, \textit{loc}, \textit{type} & Dati validi, ruolo \texttt{operator} & \texttt{DB} & La richiesta ritorna lo status \texttt{201}, edificio creato & La richiesta ritorna lo status \texttt{201}, edificio creato & RF5 \\
    \hline
    4.3 & \texttt{POST /api/v1/buildings} dati non validi & \textit{surface} neg. & Campi fuori range o mancanti & - & La richiesta ritorna con status \texttt{400} dovuto ad un errore di validazione Zod & La richiesta ritorna con status \texttt{400} dovuto ad un errore di validazione Zod & RF5, RNF5 \\
    \hline
    4.4 & \texttt{POST /api/v1/buildings} tipologia inesistente & \textit{buildingType} & Riferimento tipologia assente & \texttt{DB} & Errore \texttt{404} tipologia non trovata & Errore \texttt{404} tipologia non trovata & RF5 \\
    \hline
    4.5 & \texttt{PATCH /api/v1/buildings/:id} modifica edificio & \texttt{id}, \textit{name} & Edificio esistente & \texttt{DB} & La richiesta ritorna lo status \texttt{200}, documento aggiornato & La richiesta ritorna lo status \texttt{200}, documento aggiornato & RF5 \\
    \hline
    4.6 & \texttt{GET /api/v1/buildings/:id} edificio inesistente & \texttt{id} & ID non presente nel DB & \texttt{DB} & Errore \texttt{404} non trovato & Errore \texttt{404} non trovato & RF5 \\
    \hline
    4.7 & \texttt{PATCH /api/v1/buildings/:id} soglie COP valide & \textit{efficiencyThresholds} & Edificio esistente, \texttt{minCop} presente & \texttt{DB} & La richiesta ritorna lo status \texttt{200}, soglie salvate & La richiesta ritorna lo status \texttt{200}, soglie salvate & RF6 \\
    \hline
    4.8 & \texttt{PATCH /api/v1/buildings/:id} soglia senza \texttt{minCop} & \textit{efficiencyThresholds} & \texttt{enabled} senza \texttt{minCop} & \texttt{DB} & La richiesta ritorna con status \texttt{400} dovuto ad un errore di validazione & La richiesta ritorna con status \texttt{400} dovuto ad un errore di validazione & RF6, RNF5 \\
    \hline
    4.9 & \texttt{PATCH /api/v1/buildings/:id} soglie su teleriscaldamento & \textit{efficiencyThresholds} & Edificio a teleriscaldamento & \texttt{DB} & Errore \texttt{400} non applicabile & Errore \texttt{400} non applicabile & RF6 \\
    \hline
    4.10 & \texttt{DELETE /api/v1/buildings/:id} cancellazione a cascata & \texttt{id} & Edificio con sensori, allarmi e letture & \texttt{DB} & La richiesta ritorna lo status \texttt{200}, entit\'a collegate rimosse & La richiesta ritorna lo status \texttt{200}, entit\'a collegate rimosse & RF5 \\
    \hline
    4.11 & \texttt{GET /api/v1/buildings} ricerca per nome & \textit{name} & Edifici presenti nel DB & \texttt{DB} & La richiesta ritorna lo status \texttt{200}, lista filtrata & La richiesta ritorna lo status \texttt{200}, lista filtrata & RF3 \\
    \hline
    4.12 & \texttt{GET /api/v1/buildings} ricerca per indirizzo & \textit{address} & Edifici presenti nel DB & \texttt{DB} & La richiesta ritorna lo status \texttt{200}, lista filtrata & La richiesta ritorna lo status \texttt{200}, lista filtrata & RF3 \\
    \hline
    4.13 & \texttt{GET /api/v1/buildings} filtro per zona & \textit{zone} & Edifici presenti nel DB & \texttt{DB} & La richiesta ritorna lo status \texttt{200}, lista filtrata & La richiesta ritorna lo status \texttt{200}, lista filtrata & RF3 \\
    \hline
    4.14 & \texttt{GET /api/v1/buildings} filtro per tipologia & \textit{type} & Edifici presenti nel DB & \texttt{DB} & La richiesta ritorna lo status \texttt{200}, lista filtrata & La richiesta ritorna lo status \texttt{200}, lista filtrata & RF3 \\
    \hline
    4.15 & \texttt{GET /api/v1/buildings} filtri combinati & \textit{name}, \textit{zone}, \textit{status} & Edifici presenti nel DB & \texttt{DB} & La richiesta ritorna lo status \texttt{200}, lista filtrata & La richiesta ritorna lo status \texttt{200}, lista filtrata & RF3 \\
    \hline
    4.16 & \texttt{GET /api/v1/buildings} campi obbligatori in elenco & - & Edifici presenti nel DB & \texttt{DB} & La richiesta ritorna lo status \texttt{200}, nome, indirizzo, superficie, impianto, stato, aggiornamento & La richiesta ritorna lo status \texttt{200}, nome, indirizzo, superficie, impianto, stato, aggiornamento & RF3 \\
    \hline
    4.17 & \texttt{GET /api/v1/buildings/:id} dettaglio edificio & \texttt{id} & Edificio esistente & \texttt{DB} & La richiesta ritorna lo status \texttt{200}, anagrafica completa & La richiesta ritorna lo status \texttt{200}, anagrafica completa & RF4 \\
    \hline
    4.18 & \texttt{GET /api/v1/buildings/:id/history} serie storica & \textit{startDate}, \textit{endDate} & Letture nell'intervallo & \texttt{DB} & La richiesta ritorna lo status \texttt{200}, rilevazioni ordinate & La richiesta ritorna lo status \texttt{200}, rilevazioni ordinate & RF4, RF8 \\
    \hline
    4.19 & \texttt{GET /api/v1/buildings/:id/history} filtro per data & \textit{startDate}, \textit{endDate} & Letture nell'intervallo & \texttt{DB} & La richiesta ritorna lo status \texttt{200}, solo rilevazioni nel periodo & La richiesta ritorna lo status \texttt{200}, solo rilevazioni nel periodo & RF4, RF8 \\
    \hline
    4.20 & \texttt{GET /api/v1/buildings/:id/history} tipi di intervallo & \textit{interval} & Letture presenti & \texttt{DB} & La richiesta ritorna lo status \texttt{200}, aggregazione per intervallo & La richiesta ritorna lo status \texttt{200}, aggregazione per intervallo & RF4 \\
    \hline
    4.21 & \texttt{GET /api/v1/buildings/:id/efficiency} edificio inesistente & \texttt{id} & ID non presente nel DB & \texttt{DB} & Errore \texttt{404} non trovato & Errore \texttt{404} non trovato & RF12 \\
    \hline
    4.22 & \texttt{GET /api/v1/buildings/:id/efficiency} data non ISO & \textit{startDate} & Data malformata & - & La richiesta ritorna con status \texttt{400} dovuto ad un errore di validazione & La richiesta ritorna con status \texttt{400} dovuto ad un errore di validazione & RF12, RNF5\\
    \hline
    4.23 & \texttt{GET /api/v1/buildings/:id/efficiency} senza dati & \texttt{id}, date & Nessuna lettura nel periodo & \texttt{DB} & La richiesta ritorna lo status \texttt{200}, metriche a \texttt{null} & La richiesta ritorna lo status \texttt{200}, metriche a \texttt{null} & RF12 \\
    \hline
    4.24 & \texttt{GET /api/v1/buildings/:id/efficiency} calcolo COP & \texttt{id}, date & Consumi e temperature presenti & \texttt{DB} & La richiesta ritorna lo status \texttt{200}, COP e dispersione calcolati & La richiesta ritorna lo status \texttt{200}, COP e dispersione calcolati & RF12 \\
    \hline
    4.25 & \texttt{GET /api/v1/buildings/:id/efficiency} temperature meteo & \texttt{id}, date & Senza sonda esterna, API meteo disponibile & \texttt{DB}, meteo & La richiesta ritorna lo status \texttt{200}, COP e H da fallback meteo & La richiesta ritorna lo status \texttt{200}, COP e H da fallback meteo & RNF8 -- Fonte meteo \\
    \hline
    4.26 & \texttt{GET /api/v1/buildings/:id/efficiency} consumi gas & \texttt{id}, date & Letture gas cumulative & \texttt{DB} & La richiesta ritorna lo status \texttt{200}, consumo da m$^3 \times$ PCI & La richiesta ritorna lo status \texttt{200}, consumo da m$^3 \times$ PCI & RF12 \\
    \hline
    4.27 & \texttt{GET /api/v1/buildings/:id/efficiency} teleriscaldamento & \texttt{id}, date & Impianto a teleriscaldamento & \texttt{DB} & La richiesta ritorna lo status \texttt{200}, COP a \texttt{null} & La richiesta ritorna lo status \texttt{200}, COP a \texttt{null} & RF6, RF12 \\
    \hline
    4.28 & \texttt{GET /api/v1/buildings/:id/efficiency} calcolo efficienza & \texttt{id}, date & Consumi e temperature presenti & \texttt{DB} & La richiesta ritorna lo status \texttt{200}, metriche di efficienza & La richiesta ritorna lo status \texttt{200}, metriche di efficienza & RF12 \\
    \hline
    4.29 & \texttt{GET /api/v1/buildings/:id/efficiency} date non valide & \textit{startDate}, \textit{endDate} & Intervallo malformato & - & La richiesta ritorna con status \texttt{400} dovuto ad un errore di validazione & La richiesta ritorna con status \texttt{400} dovuto ad un errore di validazione & RF12, RNF5 \\
    \hline
    5.1 & \texttt{POST /api/v1/sensors} registrazione sensore & \textit{building}, \textit{sensorType}, \textit{location} & Edificio esistente, ruolo \texttt{admin} & \texttt{DB} & La richiesta ritorna lo status \texttt{201}, sensore \texttt{active} & La richiesta ritorna lo status \texttt{201}, sensore \texttt{active} & RF5 \\
    \hline
    5.2 & \texttt{POST /api/v1/sensors} registrazione con operator & \textit{building}, \textit{sensorType}, \textit{location} & Edificio esistente, ruolo \texttt{operator} & \texttt{DB} & La richiesta ritorna lo status \texttt{201}, sensore creato & La richiesta ritorna lo status \texttt{201}, sensore creato & RF5\\
    \hline
    5.3 & \texttt{POST /api/v1/sensors} matricola duplicata & \textit{serialNumber} & Matricola gi\'a registrata & \texttt{DB} & Errore \texttt{400} violazione unicit\'a & Errore \texttt{400} violazione unicit\'a & RF5 \\
    \hline
    5.4 & \texttt{POST /api/v1/sensors} dati non validi & \textit{sensorType}, \textit{location} & Tipo fuori enum o posizione vuota & \texttt{DB} & La richiesta ritorna con status \texttt{400} dovuto ad un errore di validazione & La richiesta ritorna con status \texttt{400} dovuto ad un errore di validazione & RF5, RNF5 \\
    \hline
    5.5 & \texttt{GET /api/v1/sensors/:id} dettaglio sensore & \texttt{id} & Sensore esistente & \texttt{DB} & La richiesta ritorna lo status \texttt{200}, dettaglio completo & La richiesta ritorna lo status \texttt{200}, dettaglio completo & RF7 \\
    \hline
    5.6 & \texttt{GET /api/v1/sensors/:id} sensore inesistente & \texttt{id} & ID non presente nel DB & \texttt{DB} & Errore \texttt{404} non trovato & Errore \texttt{404} non trovato & RF7 \\
    \hline
    5.7 & \texttt{PATCH /api/v1/sensors/:id} soglie min/max & \texttt{id}, \textit{minThreshold}, \textit{maxThreshold} & Sensore esistente & \texttt{DB} & La richiesta ritorna lo status \texttt{200}, soglie aggiornate & La richiesta ritorna lo status \texttt{200}, soglie aggiornate & RF9 \\
    \hline
    5.8 & \texttt{PATCH /api/v1/sensors/:id} sensore inesistente & \texttt{id} & ID non presente nel DB & \texttt{DB} & Errore \texttt{404} non trovato & Errore \texttt{404} non trovato & RF9 \\
    \hline
    5.9 & \texttt{DELETE /api/v1/sensors/:id} cancellazione a cascata & \texttt{id} & Sensore con letture e allarmi & \texttt{DB} & La richiesta ritorna lo status \texttt{200}, dati associati rimossi & La richiesta ritorna lo status \texttt{200}, dati associati rimossi & RF5 \\
    \hline
    5.10 & \texttt{DELETE /api/v1/sensors/:id} sensore inesistente & \texttt{id} & ID non presente nel DB & \texttt{DB} & Errore \texttt{404} non trovato & Errore \texttt{404} non trovato & RF5 \\
    \hline
    5.11 & Controllo inattivit\'a sensore oltre $2\times$ intervallo & \texttt{id} & Nessuna lettura da oltre il doppio dell'intervallo & \texttt{DB} & Stato \texttt{inactive}, \texttt{isOffline=true} & Stato \texttt{inactive}, \texttt{isOffline=true} & RF7, RNF4 \\
    \hline
    5.12 & Controllo attivit\'a sensore entro $2\times$ intervallo & \texttt{id} & Lettura recente entro il doppio dell'intervallo & \texttt{DB} & Stato \texttt{active} mantenuto & Stato \texttt{active} mantenuto & RF7 \\
    \hline
    5.13 & Sensori in manutenzione o errore con letture vecchie & \texttt{id} & Stato \texttt{maintenance} o \texttt{error} & \texttt{DB} & Stato invariato & Stato invariato & RF7 \\
    \hline
    5.14 & Sensore attivo senza alcuna lettura & \texttt{id} & \texttt{lastReading} assente & \texttt{DB} & Stato \texttt{inactive} & Stato \texttt{inactive} & RF7, RNF4 \\
    \hline
    5.15 & \texttt{GET /api/v1/sensors/:id/readings} storico letture & \texttt{id}, date & Letture presenti & \texttt{DB} & La richiesta ritorna lo status \texttt{200}, serie storica & La richiesta ritorna lo status \texttt{200}, serie storica & RF8 \\
    \hline
    5.16 & \texttt{GET /api/v1/sensors/:id/readings} sensore inesistente & \texttt{id} & ID non presente nel DB & \texttt{DB} & Errore \texttt{404} non trovato & Errore \texttt{404} non trovato & RF8 \\
    \hline
    6.1 & \texttt{GET /api/v1/alerts} elenco allarmi & \textit{status}, \textit{severity}, \textit{buildingId} & Allarmi presenti nel DB & \texttt{DB} & La richiesta ritorna lo status \texttt{200}, elenco per gravit\'a/data & La richiesta ritorna lo status \texttt{200}, elenco per gravit\'a/data & RF6 \\
    \hline
    6.2 & \texttt{PATCH /api/v1/alerts/:id} presa in carico & \texttt{id} & Allarme in stato \texttt{active} & \texttt{DB} & La richiesta ritorna lo status \texttt{200}, operatore e timestamp registrati & La richiesta ritorna lo status \texttt{200}, operatore e timestamp registrati & RF6 \\
    \hline
    6.3 & \texttt{PATCH /api/v1/alerts/:id} presa in carico non valida & \texttt{id} & Allarme gi\'a gestito o risolto & \texttt{DB} & Errore \texttt{409} transizione vietata & Errore \texttt{409} transizione vietata & RF6 \\
    \hline
    6.4 & \texttt{PATCH /api/v1/alerts/:id} risoluzione allarme & \texttt{id} & Allarme non risolto & \texttt{DB} & La richiesta ritorna lo status \texttt{200}, allarme \texttt{resolved} & La richiesta ritorna lo status \texttt{200}, allarme \texttt{resolved} & RF6 \\
    \hline
    6.5 & \texttt{GET /api/v1/alerts/:id} dettaglio allarme & \texttt{id} & Allarme esistente & \texttt{DB} & La richiesta ritorna lo status \texttt{200}, edificio, valori e link & La richiesta ritorna lo status \texttt{200}, edificio, valori e link & RF6 \\
    \hline
    6.6 & \texttt{GET /api/v1/alerts/:id} allarme inesistente & \texttt{id} & ID non presente nel DB & \texttt{DB} & Errore \texttt{404} non trovato & Errore \texttt{404} non trovato & RF6 \\
    \hline
    7.1 & \texttt{GET /api/v1/readings} intervallo date invalido & \textit{startDate}, \textit{endDate} & Date mancanti o invertite & \texttt{DB} & Errore \texttt{400}/\texttt{422} & Errore \texttt{400}/\texttt{422} & RF10, RNF5 \\
    \hline
    7.2 & \texttt{GET /api/v1/readings} esportazione CSV & \textit{buildingIds}, date & Letture nell'intervallo & \texttt{DB} & La richiesta ritorna lo status \texttt{200}, file \texttt{text/csv} & La richiesta ritorna lo status \texttt{200}, file \texttt{text/csv} & RF10 \\
    \hline
    7.3 & \texttt{GET /api/v1/readings} edificio inesistente & \textit{buildingIds} & ID non presente nel DB & \texttt{DB} & Errore \texttt{404} non trovato & Errore \texttt{404} non trovato & RF10 \\
    \hline
    7.4 & \texttt{GET /api/v1/reports} intervallo date invalido & \textit{startDate}, \textit{endDate} & Date mancanti o invertite & \texttt{DB} & Errore \texttt{400}/\texttt{422} & Errore \texttt{400}/\texttt{422} & RF10, RNF5 \\
    \hline
    7.5 & \texttt{GET /api/v1/reports} formato non valido & \textit{format} & Formato fuori enum & - & La richiesta ritorna con status \texttt{400} dovuto ad un errore di validazione & La richiesta ritorna con status \texttt{400} dovuto ad un errore di validazione & RF10, RNF5 \\
    \hline
    7.6 & \texttt{GET /api/v1/reports} edificio inesistente & \textit{buildingIds} & ID non presente nel DB & \texttt{DB} & Errore \texttt{404} non trovato & Errore \texttt{404} non trovato & RF10 \\
    \hline
    7.7 & \texttt{GET /api/v1/reports} report PDF & \textit{buildingIds}, date, \textit{pdf} & Edifici e letture presenti & \texttt{DB} & La richiesta ritorna lo status \texttt{200}, \texttt{application/pdf} & La richiesta ritorna lo status \texttt{200}, \texttt{application/pdf} & RF10 \\
    \hline
    7.8 & \texttt{GET /api/v1/reports} report Excel & \textit{buildingIds}, date, \textit{xlsx} & Edifici e letture presenti & \texttt{DB} & La richiesta ritorna lo status \texttt{200}, foglio OpenXML & La richiesta ritorna lo status \texttt{200}, foglio OpenXML & RF10 \\
    \hline
    7.9 & \texttt{GET /api/v1/reports} etichette localizzate & \textit{lang} & Lingua richiesta valida & \texttt{DB} & La richiesta ritorna lo status \texttt{200}, report nella lingua & La richiesta ritorna lo status \texttt{200}, report nella lingua & RF0 \\
    \hline
    7.10 & \texttt{GET /api/health} stato servizio & - & Nessuna & - & La richiesta ritorna lo status \texttt{200}, \texttt{\{status: "ok"\}} & La richiesta ritorna lo status \texttt{200}, \texttt{\{status: "ok"\}} & RNF4 \\
    \hline
    7.11 & \texttt{GET /api/v1/metrics} contatori aggregati & - & Sensori e allarmi presenti & \texttt{DB} & La richiesta ritorna lo status \texttt{200}, totali, attivi e consumi & La richiesta ritorna lo status \texttt{200}, totali, attivi e consumi & RF2 \\
    \hline
    7.12 & \texttt{GET /api/v1/metrics/timeseries} serie aggregate & \textit{startDate}, \textit{endDate}, \textit{interval} & Intervallo valido & \texttt{DB} & La richiesta ritorna lo status \texttt{200}, serie per periodo & La richiesta ritorna lo status \texttt{200}, serie per periodo & RF2 \\
    \hline
    7.13 & \texttt{GET /api/v1/metrics/timeseries} intervallo invalido & \textit{startDate}, \textit{endDate} & Date invertite & - & Errore \texttt{422} range invalido & Errore \texttt{422} range invalido & RF2, RNF5 -- Convalida \\
    \hline
    7.14 & \texttt{GET /api/settings} configurazione sistema & - & Ruolo \texttt{admin} & \texttt{DB} & La richiesta ritorna lo status \texttt{200}, configurazione globale & La richiesta ritorna lo status \texttt{200}, configurazione globale & RNF4 \\
    \hline
    7.15 & \texttt{PATCH /api/settings} retention e \texttt{collMod} & \textit{dataRetentionDays} & Valore entro i limiti & \texttt{DB} & La richiesta ritorna lo status \texttt{200}, retention e TTL aggiornati & La richiesta ritorna lo status \texttt{200}, retention e TTL aggiornati & RNF4 \\
    \hline
    8.1 & Ingestione lettura e \texttt{lastReading} & \textit{sensorId}, \textit{value} & Sensore esistente & \texttt{DB} & Lettura persistita, sensore aggiornato & Lettura persistita, sensore aggiornato & RF8 \\
    \hline
    8.2 & Riattivazione sensore inattivo alla nuova lettura & \textit{sensorId} & Sensore \texttt{inactive} & \texttt{DB} & Stato \texttt{active} ripristinato & Stato \texttt{active} ripristinato & RF7, RNF4 \\
    \hline
    8.3 & Lettura entro le soglie senza allarmi & \textit{value} & Valore tra min e max & \texttt{DB} & Nessun allarme creato & Nessun allarme creato & RF9 \\
    \hline
    8.4 & Superamento soglia massima & \textit{value} $>$ \textit{max} & Soglia configurata & \texttt{DB} & Allarme max attivo creato & Allarme max attivo creato & RF6 \\
    \hline
    8.5 & Superamento soglia minima & \textit{value} $<$ \textit{min} & Soglia configurata & \texttt{DB} & Allarme min attivo creato & Allarme min attivo creato & RF6 \\
    \hline
    8.6 & Lettura per sensore inesistente & \textit{sensorId} & Sensore assente nel DB & \texttt{DB} & Errore \texttt{SensorNotFound} & Errore \texttt{SensorNotFound} & RF8 \\
    \hline
    8.7 & Allarme COP sotto soglia & COP $<$ \texttt{minCop} & Soglie efficienza abilitate & \texttt{DB} & Allarme \texttt{efficiency\_cop\_low} creato & Allarme \texttt{efficiency\_cop\_low} creato & RF6 \\
    \hline
    8.8 & Deduplicazione allarmi COP attivi & COP $<$ \texttt{minCop} & Allarme attivo esistente & \texttt{DB} & Nessun duplicato creato & Nessun duplicato creato & RF6 \\
    \hline
    8.9 & Risoluzione allarme COP al recupero & COP $>$ \texttt{minCop} & Allarme attivo esistente & \texttt{DB} & Allarme risolto & Allarme risolto & RF6 \\
    \hline
    8.10 & Email di allarme multilingua agli attivi & Allarme creato & Utenti attivi, email abilitate & \texttt{DB}, \texttt{Resend} & Email nella lingua di ognuno, disabili esclusi & Email nella lingua di ognuno, disabili esclusi & RF9, RF0 \\
    \hline
    8.11 & Nessuna email con notifiche disabilitate & Allarme creato & Toggle email spento & \texttt{DB} & Nessun invio & Nessun invio & RF9  \\
    \hline
    8.12 & Destinatario fallito non blocca gli altri & Allarme creato & Un invio fallisce & \texttt{DB}, \texttt{Resend} & Gli altri ricevono l'email & Gli altri ricevono l'email & RF9 \\
    \hline
    8.13 & Simulatore emette una lettura per sensore & Sensori registrati & Simulatore avviato & \texttt{DB} & Letture valide per ogni sensore & Letture valide per ogni sensore & RF8 \\
    \hline
    8.14 & Simulatore emette consumi kW pompa di calore & Pompa di calore accesa & Simulatore avviato & \texttt{DB} & Lettura \texttt{energy\_meter} in kW & Lettura \texttt{energy\_meter} in kW & RF8 \\
    \hline
    8.15 & Simulatore emette contatore gas con resume & Caldaia a gas & Lettura precedente presente & \texttt{DB} & Lettura \texttt{gas\_meter} in m$^3$ dal valore salvato & Lettura \texttt{gas\_meter} in m$^3$ dal valore salvato & RF12  \\
    \hline
    8.16 & Simulatore incrementa odometro gas & Caldaia accesa & Simulatore avviato & \texttt{DB} & Contatore cumulativo crescente & Contatore cumulativo crescente & RF12 -- Consumi \\
    \hline
    8.17 & Messaggio MQTT valido su \texttt{sensors/+/readings} & \textit{topic}, \textit{value}, \textit{unit} & Sensore esistente, payload conforme & \texttt{DB}, \texttt{MQTT} & Lettura delegata al servizio & Lettura delegata al servizio & RF8  \\
    \hline
    8.18 & Messaggio MQTT con topic non valido & \textit{topic} & Topic senza ID sensore & \texttt{MQTT} & Messaggio ignorato senza errori & Messaggio ignorato senza errori & RNF4 \\
    \hline
    8.19 & Messaggio MQTT con JSON malformato & payload & Payload corrotto & \texttt{MQTT} & Messaggio scartato senza crash & Messaggio scartato senza crash & RNF4  \\
    \hline
    8.20 & Messaggio MQTT con payload non valido & \textit{value} & Validazione fallita & \texttt{MQTT} & Messaggio scartato senza crash & Messaggio scartato senza crash & RNF5 \\
    \hline
    9.1 & Formati email validi accettati & \textit{email} & Formati conformi & - & Nessun errore di formato & Nessun errore di formato & RNF5 \\
    \hline
    9.2 & Formati email non validi rifiutati & \textit{email} & Formati errati & - & Errore \texttt{400} formato & Errore \texttt{400} formato & RNF5 \\
    \hline
    9.3 & JSON malformato rifiutato & \textit{body} & Corpo non JSON & - & Errore \texttt{400} parsing & Errore \texttt{400} parsing & RNF5 \\
    \hline
    9.4 & Corpo richiesta vuoto rifiutato & - & Corpo assente & - & La richiesta ritorna con status \texttt{400} dovuto ad un errore di validazione & La richiesta ritorna con status \texttt{400} dovuto ad un errore di validazione & RNF5 \\
    \hline
    9.5 & Tipi di campo errati rifiutati & \textit{email}, \textit{password} & Tipi non stringa & - & La richiesta ritorna con status \texttt{400} dovuto ad un errore di validazione & La richiesta ritorna con status \texttt{400} dovuto ad un errore di validazione & RNF5  \\
    \hline
    9.6 & Campo email mancante rifiutato & - & \textit{email} assente & - & La richiesta ritorna con status \texttt{400} dovuto ad un errore di validazione & La richiesta ritorna con status \texttt{400} dovuto ad un errore di validazione & RNF5 \\
    \hline
    9.7 & Parametri token invito validi & \textit{token} & Token ben formato & \texttt{DB} & Lookup eseguita & Lookup eseguita & RF11  \\
    \hline
    9.8 & Token invito invalido rifiutato & \textit{token} & Token inesistente & \texttt{DB} & Errore \texttt{404}/\texttt{422} & Errore \texttt{404}/\texttt{422} & RF11  \\
    \hline
    9.9 & Ruoli validi accettati & \textit{role} & \texttt{admin}$|$\texttt{operator} & \texttt{DB} & Nessun errore di ruolo & Nessun errore di ruolo & RF11  \\
    \hline
    9.10 & Ruoli non validi rifiutati & \textit{role} & Ruolo fuori enum & - & La richiesta ritorna con status \texttt{400} dovuto ad un errore di validazione & La richiesta ritorna con status \texttt{400} dovuto ad un errore di validazione & RF11, RNF5 \\
    \hline
    9.11 & Campo ruolo mancante rifiutato & - & \textit{role} assente & - & La richiesta ritorna con status \texttt{400} dovuto ad un errore di validazione & La richiesta ritorna con status \texttt{400} dovuto ad un errore di validazione & RF11, RNF5 \\
    \hline
\end{xltabular}
\endgroup